% Options for packages loaded elsewhere
\PassOptionsToPackage{unicode}{hyperref}
\PassOptionsToPackage{hyphens}{url}
\PassOptionsToPackage{dvipsnames,svgnames,x11names}{xcolor}
%
\documentclass[
  11pt,
  a4paper,
]{ltjsarticle}
\usepackage{amsmath,amssymb}
\usepackage{iftex}
\ifPDFTeX
  \usepackage[T1]{fontenc}
  \usepackage[utf8]{inputenc}
  \usepackage{textcomp} % provide euro and other symbols
\else % if luatex or xetex
  \usepackage{unicode-math} % this also loads fontspec
  \defaultfontfeatures{Scale=MatchLowercase}
  \defaultfontfeatures[\rmfamily]{Ligatures=TeX,Scale=1}
\fi
\usepackage{lmodern}
\ifPDFTeX\else
  % xetex/luatex font selection
\fi
% Use upquote if available, for straight quotes in verbatim environments
\IfFileExists{upquote.sty}{\usepackage{upquote}}{}
\IfFileExists{microtype.sty}{% use microtype if available
  \usepackage[]{microtype}
  \UseMicrotypeSet[protrusion]{basicmath} % disable protrusion for tt fonts
}{}
\makeatletter
\@ifundefined{KOMAClassName}{% if non-KOMA class
  \IfFileExists{parskip.sty}{%
    \usepackage{parskip}
  }{% else
    \setlength{\parindent}{0pt}
    \setlength{\parskip}{6pt plus 2pt minus 1pt}}
}{% if KOMA class
  \KOMAoptions{parskip=half}}
\makeatother
\usepackage{xcolor}
\usepackage[margin=24mm]{geometry}
\setlength{\emergencystretch}{3em} % prevent overfull lines
\providecommand{\tightlist}{%
  \setlength{\itemsep}{0pt}\setlength{\parskip}{0pt}}
\setcounter{secnumdepth}{5}
\ifLuaTeX
  \usepackage{selnolig}  % disable illegal ligatures
\fi
\IfFileExists{bookmark.sty}{\usepackage{bookmark}}{\usepackage{hyperref}}
\IfFileExists{xurl.sty}{\usepackage{xurl}}{} % add URL line breaks if available
\urlstyle{same}
\hypersetup{
  colorlinks=true,
  linkcolor={blue},
  filecolor={Maroon},
  citecolor={Blue},
  urlcolor={blue},
  pdfcreator={LaTeX via pandoc}}

\author{}
\date{}

\begin{document}

\hypertarget{ux89e3ux7b54-02-ux884cux5217ux3068ux7ddaux5f62ux5199ux50cf}{%
\section{解答 02 ---
行列と線形写像}\label{ux89e3ux7b54-02-ux884cux5217ux3068ux7ddaux5f62ux5199ux50cf}}

\hypertarget{section}{%
\subsection{1}\label{section}}

\[
Ax=2\begin{bmatrix}1\\3\\5\end{bmatrix}-\begin{bmatrix}2\\4\\6\end{bmatrix}
=\begin{bmatrix}0\\2\\4\end{bmatrix}.
\]

行列積は列の一次結合として理解できます。

\hypertarget{section-1}{%
\subsection{2}\label{section-1}}

方程式は \texttt{x+2y+3z=0} 一本だけです。

\[
x=-2y-3z
\]

なので

\[
(x,y,z)^T=y(-2,1,0)^T+z(-3,0,1)^T.
\]

核の次元は2です。

\hypertarget{section-2}{%
\subsection{3}\label{section-2}}

第2行は第1行の2倍なので rank は1です。入力次元3に対し

\[
\dim\ker A+\operatorname{rank}A=2+1=3
\]

で次元定理を満たします。

\hypertarget{section-3}{%
\subsection{4}\label{section-3}}

\texttt{(AB)x=A(Bx)} なので \texttt{B}
が先です。写像の順序を交換すると中間状態が変わるため、一般には結果も変わります。

\hypertarget{section-4}{%
\subsection{5}\label{section-4}}

非零 \texttt{z∈ker(A)} を取ります。任意の \texttt{x} について

\[
A(x+z)=Ax+Az=Ax.
\]

異なる入力 \texttt{x} と \texttt{x+z}
が同じ出力になるため、一意復元できません。

\end{document}
